% ---------------------------------------------------------------
%  How the fourteen predictive studies of Table 2.1 test their models.
%  Counts are the verified tally in Section 2.2.6: 2 + 5 + 3 + 4 = 14.
%  Horizontal bars rather than one stacked bar, so each label has a
%  row to itself and nothing can collide.
%  Pure TikZ - needs no package beyond what main.tex already loads.
% ---------------------------------------------------------------
\begin{tikzpicture}[x=1.5cm, y=1.05cm]

% ---- one row per evaluation design: y / count / fill / edge / label ----
% One series, so one colour. Only the row the caption singles out is
% highlighted; the other three are neutral.
\foreach \y/\n/\f/\e/\lab in {%
  0/2/redfill/rededge/{Partition by population structure},
  -1/5/grayfill/neutralc/{Cross-validation or holdout, plus an independent external collection},
  -2/3/grayfill/neutralc/{Cross-validation alone},
  -3/4/grayfill/neutralc/{Random or stratified holdout, no external set}}{
    \fill[\f] (0,\y-0.28) rectangle (\n,\y+0.28);
    \draw[\e,line width=0.9pt] (0,\y-0.28) rectangle (\n,\y+0.28);
    % count, just past the end of the bar
    \node[anchor=west,font=\large\bfseries,text=black] at (\n+0.12,\y) {\n};
    % label, right-aligned in its own column so rows never touch
    \node[anchor=east,font=\small,text=black,align=right,text width=6.2cm]
      at (-0.20,\y) {\lab};
}

% ---- baseline ----
\draw[black,line width=0.9pt] (0,0.40) -- (0,-3.40);

% ---- brace marking the twelve that do not hold a lineage out by design ----
\draw[black!55,line width=0.9pt]
  (5.75,-0.72) -- (6.00,-0.72) -- (6.00,-3.28) -- (5.75,-3.28);
\node[anchor=west,font=\small,text=black,align=left,text width=3.5cm]
  at (6.12,-2.00) {twelve of the fourteen do not hold a lineage out as their
  primary design {\footnotesize(two of these report a supplementary
  structure-aware check)}};

\end{tikzpicture}
